\documentclass[12pt]{article}
\usepackage[a4paper,margin=1in]{geometry}
\usepackage{graphicx}
\usepackage{booktabs}
\usepackage{longtable}
\usepackage{enumitem}
\usepackage{hyperref}
\usepackage{amsmath}
\usepackage{siunitx}
\usepackage{xcolor}

\title{EE175AB Final Design Report\\Battery Management and Communication System (BMS)}
\author{Kaushik Vada \and Christian Kim \and Connor Stewart \and Nick Poon}
\date{\today}

\begin{document}
\maketitle
\tableofcontents
\newpage

\section{What the project is all about}

\subsection{Executive Summary}
This project delivers a complete \textbf{Battery Management and Communication System (BMS)} for a \textbf{10s1p lithium-ion battery pack}. The system combines:
\begin{itemize}[leftmargin=1.2cm]
    \item high-voltage battery sensing and protection,
    \item an isolated embedded control and communication layer, and
    \item a desktop dashboard for real-time monitoring, command/control, and firmware update support.
\end{itemize}

The architecture is intentionally split between a battery domain and a host-computer domain to improve safety and reduce fault propagation risk. On the embedded side, firmware on an STM32F303 platform acquires cell telemetry, measures thermistors, estimates current through a shunt path, controls cooling fan operation with automatic and manual modes, and streams telemetry over USB CDC. On the host side, a PyQt6 application embeds a web frontend to visualize live data and send commands back to the controller.

\subsection{Purpose and Application Context}
The design targets battery-powered engineering prototypes where safe operation, observability, and controlled interfacing with a development computer are required. Example use cases include:
\begin{itemize}[leftmargin=1.2cm]
    \item laboratory battery characterization,
    \item educational power electronics and controls experiments,
    \item embedded systems validation with pack-level telemetry,
    \item future integration with AI/ML-based state estimation pipelines.
\end{itemize}

\subsection{System Photo and Block-Diagram Artifacts}
The repository includes project media used for design communication and implementation traceability, including:
\begin{itemize}[leftmargin=1.2cm]
    \item \texttt{Senior Design Block Diagrams.pdf},
    \item schematic PDFs for BMS and USB firmware branches,
    \item GUI screenshots and enclosure renderings referenced in weekly reports,
    \item firmware and project details documentation in PDF/\LaTeX{} form.
\end{itemize}

\subsection{Delivered Subsystems}
\begin{enumerate}[leftmargin=1.2cm]
    \item \textbf{Embedded Firmware (USB-Connection-Firmware):} BQ76930 communication, current sensing, thermistor conversion, fan control, USB telemetry.
    \item \textbf{Embedded Firmware (Electronic Load branch):} related STM32 infrastructure for hardware-in-loop and load-control workflows.
    \item \textbf{Desktop Dashboard (BMS\_Dashboard):} PyQt6 launcher, serial parser, command bridge, frontend rendering, packaging/update workflow.
    \item \textbf{Project Documentation:} weekly reports, architecture summaries, firmware plans, and project detail reports.
\end{enumerate}

\section{Theoretical basis for the project}

\subsection{Battery Monitoring Theory}
A 10-cell series lithium-ion pack requires per-cell voltage visibility because pack-level voltage alone cannot detect cell imbalance. The design therefore centers on cell-resolved telemetry and safety checks (OV/UV style fault logic) and uses a dedicated battery monitor IC (BQ76930 family usage documented in firmware and project reports).

\subsection{Current Measurement Theory}
Pack current is estimated through a shunt resistor and coulomb-counter readout path. In firmware, conversion follows:
\[
I_{pack} = N_{LSB}\times I_{LSB},\quad I_{LSB}=\SI{0.000422}{A/LSB}
\]
where \(I_{LSB}\) is derived from BQ76930 voltage-per-LSB and a \SI{20}{m\ohm} shunt value.

\subsection{Thermistor Temperature Theory}
Temperature conversion uses a Beta-model NTC equation with parameters:
\begin{itemize}[leftmargin=1.2cm]
    \item \(R_0 = \SI{10}{k\ohm}\) at \SI{25}{\celsius},
    \item \(\beta = 3950\),
    \item divider supply near \SI{3.3}{V}.
\end{itemize}
Firmware computes divider resistance and applies a Steinhart/Beta approximation:
\[
\frac{1}{T}=\frac{1}{T_0}+\frac{1}{\beta}\ln\left(\frac{R}{R_0}\right)
\]
then reports temperature in \si{\celsius}.

\subsection{Fan Control Theory}
Thermal management uses PWM fan actuation and tachometer feedback. The firmware supports:
\begin{itemize}[leftmargin=1.2cm]
    \item automatic mode with piecewise temperature-to-duty mapping,
    \item manual duty override command path from GUI,
    \item hysteresis to suppress threshold chatter,
    \item moving-average RPM estimation and timeout-based stall detection.
\end{itemize}

\subsection{Isolation and Safety Theory}
The project architecture applies galvanic isolation between battery and PC domains so high-energy transients are less likely to couple into user-facing USB interfaces. Safety is layered:
\begin{itemize}[leftmargin=1.2cm]
    \item IC/firmware protections (OV/UV/current/thermal status),
    \item comparator and hardware disable paths in schematic plan,
    \item host-side visibility and command arbitration through structured telemetry.
\end{itemize}

\section{Design Specifications}

Design goals were defined as measurable, testable specifications compatible with EE175 expectations.

\subsection{Electrical and Battery Specifications}
\begin{longtable}{p{5cm}p{9cm}}
\toprule
\textbf{Specification} & \textbf{Target / Requirement} \\
\midrule
Cell configuration & 10s1p lithium-ion (Molicel P42A context in project docs) \\
Nominal pack voltage & Approximately \SI{36}{V} \\
Pack capacity (nominal) & \SI{4.2}{Ah} \\
Per-cell monitoring & 10 cell voltages captured and transmitted \\
Pack current estimate & Shunt-based measurement via BQ76930 current registers \\
Thermal channels & 10 NTC channels sampled through STM32 ADC path \\
\bottomrule
\end{longtable}

\subsection{Control and Communication Specifications}
\begin{longtable}{p{5cm}p{9cm}}
\toprule
\textbf{Specification} & \textbf{Target / Requirement} \\
\midrule
Host interface & USB-C / USB CDC serial telemetry \\
Default serial speed & 115200 baud (GUI and firmware workflow) \\
Fan control output & PWM control with auto and manual modes \\
Fan telemetry & Duty, auto/manual state, and RPM reporting \\
Data protocol robustness & Multi-format parser support and frame assembly resilience \\
Desktop runtime & PyQt6 launcher with embedded web frontend \\
\bottomrule
\end{longtable}

\subsection{Software/Usability Specifications}
\begin{itemize}[leftmargin=1.2cm]
    \item Auto-detect serial port workflow with exclusion rules for known non-target ports.
    \item Real-time telemetry display for cell voltages, temperatures, current, and fan metrics.
    \item Bidirectional command path (GUI to MCU) with fallback queue polling.
    \item Build/package support for Windows, macOS, and Linux release artifacts.
\end{itemize}

\subsection{Safety and Reliability Specifications}
\begin{itemize}[leftmargin=1.2cm]
    \item Over-current and thermal fault observability.
    \item Sensor conversion fault coding for disconnected/shorted/error states.
    \item Isolation boundary between battery and host domains.
    \item Reportable system status flags (e.g., SYS\_STAT, load presence, fan mode/status).
\end{itemize}

\section{Implementation Strategy}

\subsection{High-Level Architecture (System Block Design)}
The implementation strategy separated development into parallel tracks:
\begin{enumerate}[leftmargin=1.2cm]
    \item \textbf{Hardware + firmware track:} establish MCU interfaces, BMS monitor communication, thermal sensing, and fan control.
    \item \textbf{Desktop software track:} establish serial ingestion, parsing, UI rendering, and command channel.
    \item \textbf{Documentation + verification track:} weekly technical reports, firmware architecture notes, and evolving project detail documents.
\end{enumerate}

\subsection{Repository-Level Decomposition}
\begin{itemize}[leftmargin=1.2cm]
    \item \texttt{USB-Connection-Firmware/}: primary embedded BMS telemetry/control firmware.
    \item \texttt{Electronic\_Load\_Firmware/}: related STM32 firmware branch and drivers.
    \item \texttt{BMS\_Dashboard/}: desktop app, backend parser/bridge, frontend scene, packaging and updater.
    \item \texttt{Project Details/} and \texttt{Weekly Reports/}: design intent, progress, and architecture snapshots.
\end{itemize}

\subsection{Communication Strategy}
The strategy intentionally used plain serial framing and tolerant parsing to accelerate bring-up:
\begin{itemize}[leftmargin=1.2cm]
    \item firmware emits structured lines and status tokens,
    \item Python parser reconstructs coherent frames from mixed-line streams,
    \item frontend consumes normalized dictionaries/JSON-like objects for display,
    \item command bridge supports GUI-triggered control actions.
\end{itemize}

\subsection{Risk Mitigation Strategy}
\begin{itemize}[leftmargin=1.2cm]
    \item Auto-detect BQ address and CRC mode to reduce deployment fragility.
    \item Treat malformed/mixed serial logs as expected, not exceptional.
    \item Maintain fallback mock data and command queue mechanisms to continue UI testing when hardware is unavailable.
    \item Keep firmware and GUI independent enough for staged integration.
\end{itemize}

\section{Actual Implementation}

\subsection{Embedded Firmware Implementation}
\subsubsection{MCU Platform and Peripheral Bring-Up}
Implementation was built around STM32CubeIDE-generated scaffolding plus custom logic. The firmware includes initialization and runtime usage of:
\begin{itemize}[leftmargin=1.2cm]
    \item I2C for BQ76930 communication,
    \item ADC channels for NTC sensing,
    \item timers for PWM fan control and tach period capture,
    \item USB device CDC interface for host communications.
\end{itemize}

\subsubsection{Battery Monitor Communication}
Key implementation details include:
\begin{itemize}[leftmargin=1.2cm]
    \item auto-detection of BMS address candidates (\texttt{0x08}/\texttt{0x18}),
    \item optional CRC8 support with I2C transactions,
    \item register read/write wrappers for SYS\_STAT/control/protection/current paths,
    \item conversion utilities for current and voltage reporting.
\end{itemize}

\subsubsection{Thermistor Acquisition and Conversion}
The firmware configures ADC sampling with long sample times for higher-impedance thermistor networks, converts raw ADC readings to voltage, then converts voltage to temperature with explicit fault code branches for:
\begin{itemize}[leftmargin=1.2cm]
    \item hardware/timeout conversion errors,
    \item open-circuit and short-circuit sensor conditions.
\end{itemize}

\subsubsection{Fan Control and Diagnostics}
Fan control implementation includes:
\begin{itemize}[leftmargin=1.2cm]
    \item PWM duty mapping to timer ARR/CCR settings,
    \item tachometer capture for RPM calculation,
    \item moving average smoothing and timeout handling,
    \item command-driven manual override.
\end{itemize}

\subsection{Desktop Software Implementation}
\subsubsection{Runtime/Launcher Layer}
\texttt{gui\_launcher.py} provides a desktop shell with:
\begin{itemize}[leftmargin=1.2cm]
    \item PyQt6 windowing and embedded WebEngine frontend,
    \item bridge exposure (\texttt{sendCommand}) from JS to Python,
    \item serial worker thread lifecycle management,
    \item update flow integration and cross-platform icon/resource handling.
\end{itemize}

\subsubsection{Serial Parsing and Telemetry Reconstruction}
\texttt{backend/data\_stream.py} implements robust parsing of line-oriented telemetry using regex-based extraction and frame reconstruction for:
\begin{itemize}[leftmargin=1.2cm]
    \item indexed cell/temperature tokens,
    \item pack current and fan fields,
    \item mixed-format line streams,
    \item stale frame timeout finalization and last-known state propagation.
\end{itemize}

\subsubsection{Frontend and Interaction Layer}
The frontend stack (HTML/CSS/JS loaded by the launcher) renders live battery and thermal state while allowing command dispatch to firmware through the QWebChannel bridge and queue fallback path.

\subsubsection{Packaging and Distribution}
The BMS dashboard repository includes:
\begin{itemize}[leftmargin=1.2cm]
    \item dependency manifests,
    \item PyInstaller spec and release docs,
    \item helper scripts for icons and update tooling,
    \item release links for Windows/macOS/Linux binaries.
\end{itemize}

\subsection{Engineering Process Implementation}
Weekly reports track implementation increments such as fan auto-curve tuning, parser upgrades, command reliability improvements, and host runtime fixes. Project details documents preserve architecture decisions and evolving firmware strategy.

\section{Validation/Verification of Functionality}

\subsection{Verification Approach}
Validation was executed using staged testing:
\begin{enumerate}[leftmargin=1.2cm]
    \item module-level firmware checks (ADC conversion, BMS register access, PWM control),
    \item parser-level host checks on mixed serial logs,
    \item GUI integration checks for display refresh and command dispatch,
    \item end-to-end bench tests with hardware connected when available.
\end{enumerate}

\subsection{Measured/Reported Verification Items}
\begin{longtable}{p{5cm}p{4cm}p{5cm}}
\toprule
\textbf{Test Item} & \textbf{Pass Criteria} & \textbf{Observed Outcome} \\
\midrule
BQ address detection & Correctly locks to valid address & Implemented with dual-address probe and runtime selection \\
CRC compatibility & Correct read/write framing with/without CRC & Firmware supports both modes and dynamic handling \\
NTC conversion path & Stable conversion under intended sample timing & Long ADC sample time selected to improve stability \\
Fan PWM control & Duty updates reflected in fan behavior & Manual and auto duty control paths implemented \\
Fan RPM telemetry & RPM values reported and filtered & Tach capture + averaging + timeout logic implemented \\
Serial parser resilience & Reconstructs frames from mixed lines & Multi-format regex parser and pending-frame logic implemented \\
GUI command reliability & Commands reach firmware consistently & Direct bridge + polling fallback command queue present \\
\bottomrule
\end{longtable}

\subsection{Pass/Fail Conditions}
\textbf{Pass:}
\begin{itemize}[leftmargin=1.2cm]
    \item ten cell voltages and thermistor channels are visible in GUI,
    \item current and fan telemetry update in real time,
    \item manual fan command changes duty and reports status,
    \item no critical parser drops under expected mixed serial logs,
    \item protection/fault indicators can be surfaced to host software.
\end{itemize}

\textbf{Fail:}
\begin{itemize}[leftmargin=1.2cm]
    \item repeated I2C telemetry corruption with no fallback,
    \item unstable or non-physical temperature reporting under normal conditions,
    \item command path deadlock or unresponsive GUI control loop,
    \item inability to maintain data stream continuity for system state decisions.
\end{itemize}

\subsection{Current Validation Status}
Based on weekly-report milestones, the project moved from architecture and enclosure definition (Oct 2025), to firmware and telemetry bring-up (Jan 2026), to parser and GUI reliability hardening plus fan-control integration (Feb 2026), and then toward live hardware integration/calibration and fault-management expansion.

\section{Future Improvements}

\subsection{Firmware Enhancements}
\begin{itemize}[leftmargin=1.2cm]
    \item Add stronger packet framing (length + CRC) for all telemetry records.
    \item Introduce explicit firmware state machine enums with fault latching/clear policy.
    \item Add calibration storage for thermistor offsets and current-sense gain correction.
    \item Expand diagnostics counters (I2C retries, parser frame drops, watchdog resets).
\end{itemize}

\subsection{Control and Safety Enhancements}
\begin{itemize}[leftmargin=1.2cm]
    \item Implement refined thermal derating logic rather than threshold-only responses.
    \item Add SOC/SOH estimation module (coulomb count + OCV fusion) for better state reporting.
    \item Harden redundant hardware trip pathways and post-fault forensic logging.
\end{itemize}

\subsection{Desktop/UX Enhancements}
\begin{itemize}[leftmargin=1.2cm]
    \item Add persistent session logging and exportable CSV/JSON datasets.
    \item Add trend charts for voltage sag, thermal gradients, and fan response.
    \item Provide explicit alert panes with acknowledgement workflow and fault timeline.
    \item Add guided configuration pages for thresholds and protection profiles.
\end{itemize}

\subsection{Verification and Documentation Enhancements}
\begin{itemize}[leftmargin=1.2cm]
    \item Build a formal test matrix with traceability from requirements to test evidence.
    \item Automate parser regression tests with recorded serial fixture files.
    \item Add reproducible bring-up and factory-reset procedures.
    \item Expand final documentation with a reproducibility appendix (BOM revisions, wiring maps, and known issues).
\end{itemize}


\subsection{Repository-wide artifact review (files, images, PDFs, schematics)}
To satisfy complete repository traceability for the EE175 folder, a full-file manifest was generated and reviewed at:
\begin{itemize}[leftmargin=1.2cm]
    \item \texttt{Final Report/ee175\_artifact\_manifest.tsv}
\end{itemize}

This manifest enumerates every file found under \texttt{EE175 - Senior Design} with path, extension, byte count, and SHA-256 fingerprint prefix. The current manifest snapshot contains 4506 artifacts (excluding itself), spanning firmware source, generated build outputs, dashboard runtime assets, PDFs, CAD/schematic exports, and weekly-report images.

\textbf{Audit highlights used to inform this report:}
\begin{itemize}[leftmargin=1.2cm]
    \item \textbf{PDF/schematic set reviewed:} block diagrams, BMS schematics, e-load schematics, firmware architecture/plan PDFs, CubeMX project report PDF, weekly-report PDFs, and syllabus/template PDFs.
    \item \textbf{Image set reviewed:} dashboard branding/icons, weekly-report figures (BMS GUI, enclosure, battery pack, firmware screenshots), and schematic crop exports in firmware temp folders.
    \item \textbf{Firmware source reviewed:} primary STM32 application and USB CDC stack usage in \texttt{USB-Connection-Firmware/Core/Src/main.c} and related CubeMX outputs.
    \item \textbf{Dashboard source reviewed:} launcher/backend/frontend implementation for serial ingestion, command bridging, and 3D UI presentation.
\end{itemize}

\section*{Conclusion}
This EE175 project demonstrates a complete cross-domain engineering design: battery pack instrumentation, embedded control, host communication, and operator-facing visualization. The repository evidences iterative progress from concept and architecture to integrated firmware/dashboard behavior, with clear opportunities for production-grade hardening in future revisions.

\end{document}
